\section{Introduction}
\label{sec:intro}

Rice (\textit{Oryza sativa}) feeds about half of the world's population and is grown on more than 160 million
hectares \cite{faostat2024}. Pathogens and pests cause a mean yield loss of roughly 30\,\% in rice, the highest among
the five major food crops \cite{savary2019}. Rice is Pakistan's second most important cereal after wheat
\cite{gopk2024}. A foliar disease can be controlled only in a short window, and there are too few extension officers
to inspect every field, so an image-based diagnosis that works on photos a farmer can take would remove a real
bottleneck.

Automated recognition has moved from hand-made features to fine-tuned convolutional networks, transformers and generative
augmentation (Section~\ref{sec:related}), but one formulation has stayed the same: one label is predicted from one strongly
downsampled view of the photo. That is a poor match to the task. A lesion can fill well under 1\,\% of the frame, so a
photo of $1881\times1881$ pixels \cite{do2019} shrunk to $224\times224$ keeps about 1 pixel in 70, and the pale centre or
halo that separates brown spot from leaf blast is averaged away.

Three gaps follow from the literature reviewed in Section~\ref{sec:related}.
(i) Resolution is fixed early and never varied, although region aggregation is standard for gigapixel pathology images
\cite{ilse2018, lu2021clam}. (ii) Background suppression is either hand-made or needs annotation; a learned method that
uses photo-level labels only has not been evaluated for rice. (iii) Reporting practice makes numbers hard to compare:
the order of splitting and augmenting is often unstated, accuracy on imbalanced data is the usual metric, and
cross-dataset tests are rare. Testing problems of this kind are common in the field \cite{xu2024}, and a model above
90\,\% on several benchmarks has been reported at 58.15\,\% on an unseen rice set \cite{thakur2023}.

We therefore treat resolution as a variable and the protocol as part of the contribution. The method keeps lesion-scale
detail by cutting each photo into high-detail tiles at a common leaf scale, describing every tile with a frozen
self-supervised encoder, and letting a small attention head learn from the photo label which tiles matter. This is
multiple instance learning (MIL): a photo is a labelled bag of unlabelled tiles. The contributions are:
\begin{itemize}
\item a tile-and-attention pipeline for rice leaf disease that trains only about a quarter of a million parameters,
and its heat maps (Section~\ref{sec:method});
\item a protocol that groups near-duplicates and capture sessions, fixes the split before augmentation, reports
spread and two bootstrap ranges, and states a rule for ``A beats B'' (Section~\ref{sec:setup});
\item controlled comparisons on resolution, pooling, Otsu preprocessing and encoders, a zero-shot test on field photos,
a deletion test of the heat maps, a cost measurement, and a partition-order experiment on how much a faulty split
raises the score (Section~\ref{sec:results}).
\end{itemize}
Five hypotheses, fixed before the experiments, structure the study:
\textbf{H1}, tiles beat whole-photo input with the same encoder;
\textbf{H2}, the learned attention is at least as good as mean pooling, max pooling and Otsu preprocessing;
\textbf{H3}, frozen features hold up on unseen field photos;
\textbf{H4}, the heat map points at real symptoms;
\textbf{H5}, the cost is reasonable and can be reduced.
Section~\ref{sec:related} reviews related work, Section~\ref{sec:method} describes the method,
Section~\ref{sec:setup} the data and protocol, Section~\ref{sec:results} the results, and
Sections~\ref{sec:discussion}--\ref{sec:conclusion} discuss them, state the limits and conclude.
